\documentclass{article}

\usepackage{amsthm,amsfonts}
\usepackage{mathpartir}
\usepackage{mathtools}
\usepackage{tikz-cd}
\usepackage{hyperref,cleveref}
\usepackage{color,soul}
\usepackage{mathabx}

\theoremstyle{definition}
\newtheorem{definition}{Definition}[section]

% \theoremstyle{theorem}
\newtheorem{theorem}[definition]{Theorem}
\newtheorem{lemma}[definition]{Lemma}
\newtheorem{corollary}[definition]{Corollary}
\newtheorem{problem}[definition]{Problem}
\newtheorem{exercise}[definition]{Exercise}
\newtheorem{example}[definition]{Example}

\newenvironment{construction}{
\begin{proof}[Construction]}{
\end{proof}}

\newcommand{\types}{\mathcal T}
\newcommand{\type}{\ \textsc{type}}
\newcommand{\terms}{\mathcal S}
% \newcommand{\contexts}{\mathcal C}
% \newcommand{\context}{\textsc{ctx}}
\newcommand{\T}{\mathbb T}
\newcommand{\C}{\mathcal C}
\newcommand{\D}{\mathcal D}
\newcommand{\A}{\mathcal A}
\newcommand{\M}{\mathcal M}
\newcommand{\W}{\mathcal W}
\newcommand{\F}{\mathcal F}
\newcommand{\Set}{{\mathcal S}et}
\newcommand{\syncat}[1]{\C [#1]}
% \newcommand{\defeq}{\coloneqq}
% \newcommand{\interp}[1]{\lceil #1 \rceil}
\newcommand{\seq}{\doteq}
% \newcommand{\lists}{\mathcal Lists}
% \newcommand{\variables}{\mathcal Var}
\newcommand{\Epsilon}{\mathrm E}
\newcommand{\Zeta}{\mathrm Z}
\newcommand{\mor}{\mathrm {mor}}
\newcommand{\op}{\mathrm {op}}
\newcommand{\grpd}{\mathcal G}
\newcommand{\Id}{\mathtt {Id}}
\newcommand{\Path}{\mathrm {Path} \ }

\tikzset{pullbackcorner/.style={minimum size=1.2ex,path picture={
      \draw[opacity=1,black,-,#1] (-0.5ex,-0.5ex) -- (0.5ex,-0.5ex) -- (0.5ex,0.5ex);%
}}}

\newcommand{\pullback}{\arrow[dr, phantom, "" {pullbackcorner} , very near start]}

\title{HW 7 (Mastermath HoTT)}
\author{Paige Randall North}

\begin{document}

\maketitle

\noindent\textbf{Problem 1.} Pullbacks.

\begin{enumerate}
  \item[(a)] (Pullback pasting lemma.) Consider a diagram of the following form in a category $\C$.
    \[
      \begin{tikzcd}
        \bullet \ar[r] \ar[d] & \bullet \ar[r] \ar[d] & \bullet \ar[d] \\
        \bullet \ar[r] & \bullet \ar[r] & \bullet
      \end{tikzcd}
    \]
    Suppose the right square is a pullback. Show that the outer rectangle is a pullback if and only if the left square is a pullback.

  \item[(b)] (Sections from diagonals.) Consider a pullback square
    \[
      \begin{tikzcd}
        P \ar[r,"g"] \ar[d,"f"] \pullback & B \ar[d,"h"] \\
        A \ar[r,"k"] & C
      \end{tikzcd}
    \]
    in a category $\C$. Show that there is a bijection between sections of $f$ (i.e.\ morphisms $s : A \to P$ with $f s = 1_A$) and diagonals of the square (i.e.\ morphisms $d : A \to B$ with $h d = k$).

  \item[(c)] Let $\C$ be a category with a terminal object $*$ and binary products. Show that for any object $A$ of $\C$, the following square is a pullback.
    \[
      \begin{tikzcd}
        A \times A \ar[r,"\pi_1"] \ar[d,"\pi_0"] & A \ar[d,"!_A"] \\
        A \ar[r,"!_A"] & *
      \end{tikzcd}
    \]

  \item[(d)] Let $\C$ be a category with binary products. Show that for any morphism $f : X \to Y$ in $\C$, the following square is a pullback.
    \[
      \begin{tikzcd}
        X \times Y \ar[r,"f \times 1_Y"] \ar[d,"\pi_0"] & Y \times Y \ar[d,"\pi_0"] \\
        X \ar[r,"f"] & Y
      \end{tikzcd}
    \]
\end{enumerate}

\hspace{1em}

Recall the following definition.
\begin{definition}
  Given a choice of path object $\Path A$ for each object $A$, we can produce a \emph{mapping path space} factorization on the category.

  We take a morphism $f: X \to Y$ to the following factorization
  \[
    \begin{tikzcd}
      X \ar[r,"\lambda_f"] & X \times_\bullet \Path Y \ar[r,"\epsilon_f"] & Y
    \end{tikzcd}
  \]
  where $X \times_\bullet \Path Y$ is the following pullback
  \[
    \begin{tikzcd}
      X \times_\bullet \Path Y \ar[r," "] \ar[d," "] \pullback & \Path Y \ar[d,"\pi_0 \epsilon "]
      \\
      X \ar[r,"f "] & Y
    \end{tikzcd}
  \]
  The map $\epsilon_f$ is given by projecting to $\Path Y$ and then taking $\pi_1 \epsilon$. The map $\lambda_f$ is induced by $1_X$ and $r_Y f$.
\end{definition}

\noindent\textbf{Problem 2.} (The right component of the mapping path space factorization is a display map.)
Suppose that $\C$ is a clan (Def. 2.7 in the notes) with a choice of path objects for each object. Show that $\epsilon_f$ is in $\D$. (Hint: Use Problem 1 (a), (c) and (d))

\hspace{1em}

\noindent\textbf{Problem 3.} (wfs induces a clan on the fibrant objects.) Consider a weak factorization system $(\mathcal{L}, \mathcal{R})$ on a category $\C$ with finite limits. Let $\C_{\mathcal{F}}$ denote the full subcategory of $\C$ spanned by the objects $X$ such that the unique map $X \to *$ is in $\mathcal{R}$. Show that:
\begin{enumerate}
  \item[(i)] The wfs $(\mathcal{L}, \mathcal{R})$ induces a wfs $(\mathcal{L}_{\mathcal{F}}, \mathcal{R}_{\mathcal{F}})$ on $\C_{\mathcal{F}}$.
  \item[(ii)] $(\C_{\mathcal{F}}, \mathcal{R}_{\mathcal{F}})$ is a clan.
\end{enumerate}

\end{document}